\chapter{Future Scope}

Building on the successful implementation and conference publication of RE-TabSyn, several promising avenues for future research emerge. This chapter outlines near-term improvements and longer-term research directions that could extend the paradigm of controllable tabular synthesis.

\section{Architectural Improvements}
\begin{itemize}
    \item Upgrading the VAE encoder from an MLP to a full Transformer architecture to capture inter-column dependencies through self-attention, potentially reducing the fidelity gap between RE-TabSyn and TabSyn.
    \item Implementing separate noise schedules for numerical and categorical features in the latent space to reduce reconstruction errors.
    \item Scaling the Diffusion Transformer from 4 to 8--12 layers for improved capacity on high-dimensional datasets.
    \item Integrating DDIM sampling to reduce denoising steps from 1,000 to 50--100 without significant quality loss.
\end{itemize}

\section{Enhanced Privacy and Security}
\begin{itemize}
    \item Integrating formal Differential Privacy through DP-SGD during diffusion training. Preliminary prototyping with the Opacus library demonstrated feasibility at $\epsilon = 2.73$.
    \item Developing a federated RE-TabSyn framework where multiple institutions collaboratively train a global model without pooling raw data, which is particularly relevant for anti-money laundering applications.
    \item Conducting formal Membership Inference Attack evaluation using shadow model techniques to strengthen privacy claims.
\end{itemize}

\section{Extended Synthesis Capabilities}
\begin{itemize}
    \item Expanding to multi-table relational synthesis with primary and foreign key constraints, building on insights from RelDDPM \cite{liu2024relddpm}.
    \item Extending CFG beyond binary class labels to support multi-class targets, continuous regression conditioning, and arbitrary attribute constraints.
    \item Adapting the architecture for temporal tabular data such as transaction logs and stock prices, which exhibit strong sequential dependencies and regime changes.
    \item Enabling fairness-aware generation by controlling demographic attributes simultaneously with class labels to support bias auditing in lending and insurance.
\end{itemize}

\section{Evaluation and Deployment}
\begin{itemize}
    \item Testing RE-TabSyn on additional domains including healthcare, cybersecurity, and manufacturing to validate generalizability.
    \item Developing causal fidelity metrics that measure whether synthetic data preserves causal relationships, not just correlations.
    \item Building a REST API for on-demand synthetic data generation with specified minority ratios, integrated with existing MLOps pipelines.
    \item Adding a constraint enforcement layer to ensure generated data satisfies hard business rules such as valid age ranges and inter-column dependencies.
\end{itemize}
